% Шаблон курсовой работы
% Лаборатория математики ФН1
% Компиляция: pdflatex coursework.tex && pdflatex coursework.tex

\documentclass[14pt, a4paper]{extarticle}

\usepackage[T2A]{fontenc}
\usepackage[utf8]{inputenc}
\usepackage[russian]{babel}
\usepackage[left=30mm, right=15mm, top=20mm, bottom=20mm]{geometry}
\usepackage{amsmath, amssymb, amsthm}
\usepackage{graphicx}
\usepackage{setspace}
\usepackage{indentfirst}
\usepackage{hyperref}

\onehalfspacing
\setlength{\parindent}{1.25cm}

\theoremstyle{plain}
\newtheorem{theorem}{Теорема}[section]
\newtheorem{lemma}[theorem]{Лемма}
\newtheorem{corollary}[theorem]{Следствие}

\theoremstyle{definition}
\newtheorem{definition}[theorem]{Определение}
\newtheorem{example}[theorem]{Пример}

\hypersetup{colorlinks=true, linkcolor=black, citecolor=black, urlcolor=blue}

% ---------- Данные работы: заполните ----------
\newcommand{\workTitle}{Название курсовой работы}
\newcommand{\workAuthor}{Фамилия И. О.}
\newcommand{\workGroup}{ФН1-11Б}
\newcommand{\workSupervisor}{должность, Фамилия И. О.}
\newcommand{\workYear}{2026}
% ----------------------------------------------

\begin{document}

% ================= Титульный лист =================
\begin{titlepage}
  \centering
  \large

  Лаборатория математики ФН1

  \vfill

  {\Large\bfseries Курсовая работа}

  \vspace{1cm}

  {\Large \workTitle}

  \vfill

  \begin{flushright}
    \begin{minipage}{0.5\textwidth}
      Выполнил(а): студент(ка) группы \workGroup\\
      \workAuthor

      \vspace{0.7cm}

      Руководитель:\\
      \workSupervisor
    \end{minipage}
  \end{flushright}

  \vfill

  Москва, \workYear
\end{titlepage}

\setcounter{page}{2}
\tableofcontents
\newpage

% ================= Введение =================
\section{Введение}

Опишите постановку задачи, её актуальность и место среди известных результатов.
Сформулируйте цель работы и перечислите задачи, которые требуется решить.

% ================= Основная часть =================
\section{Постановка задачи}

Приведите строгую математическую формулировку. Все обозначения вводите
явно при первом употреблении.

\begin{definition}
  Текст определения.
\end{definition}

\begin{theorem}\label{th:main}
  Формулировка основного результата.
\end{theorem}

\begin{proof}
  Доказательство теоремы~\ref{th:main}.
\end{proof}

\section{Метод решения}

Опишите используемый метод и обоснуйте его применимость. Формулы, на которые
есть ссылки в тексте, нумеруйте:
\begin{equation}\label{eq:main}
  \frac{\partial u}{\partial t} = a^2 \frac{\partial^2 u}{\partial x^2}.
\end{equation}

Из~\eqref{eq:main} следует \ldots

\section{Результаты вычислительного эксперимента}

% \begin{figure}[h]
%   \centering
%   \includegraphics[width=0.8\textwidth]{plot.pdf}
%   \caption{Подпись к рисунку располагается снизу}
%   \label{fig:plot}
% \end{figure}

\begin{table}[h]
  \centering
  \caption{Подпись к таблице располагается сверху}
  \begin{tabular}{|c|c|c|}
    \hline
    Шаг сетки & Погрешность & Порядок \\
    \hline
    $0{,}1$   & $2{,}1\cdot10^{-3}$ & --- \\
    $0{,}05$  & $5{,}3\cdot10^{-4}$ & $1{,}99$ \\
    \hline
  \end{tabular}
\end{table}

% ================= Заключение =================
\section{Заключение}

Перечислите полученные результаты и соотнесите их с задачами из введения.

% ================= Литература =================
\begin{thebibliography}{9}
  \bibitem{fikhtengolts}
    \emph{Фихтенгольц Г.\,М.} Курс дифференциального и интегрального
    исчисления. --- М.: Физматлит, 2003.
  \bibitem{samarskii}
    \emph{Самарский А.\,А.} Теория разностных схем. --- М.: Наука, 1989.
\end{thebibliography}

\end{document}
